% Notes de présentation (orateur) — Autoencodeurs et débruitage d'images
% Ce document est conçu pour VOUS (présentateurs) : explications détaillées, notations,
% et "script" slide par slide, avec décryptage des formules.
%
% Compilation (dans /docs) : pdflatex -interaction=nonstopmode -halt-on-error notes_presentateur.tex

\documentclass[11pt,a4paper]{article}

\usepackage[T1]{fontenc}
\usepackage[utf8]{inputenc}
\usepackage[french]{babel}

\usepackage{amsmath,amssymb,mathtools}
\usepackage{geometry}
\geometry{margin=2.3cm}

\usepackage{xcolor}
\usepackage{enumitem}
\setlist[itemize]{leftmargin=1.2em}
\setlist[enumerate]{leftmargin=1.2em}

\usepackage{xurl}
\usepackage{hyperref}
\hypersetup{colorlinks=true, linkcolor=blue!50!black, urlcolor=blue!50!black, citecolor=blue!50!black}
\urlstyle{same}

\newcommand{\R}{\mathbb{R}}
\newcommand{\E}{\mathbb{E}}
\newcommand{\norm}[1]{\left\lVert #1 \right\rVert}

\title{Notes de présentation (explications détaillées)\\Autoencodeurs et débruitage d'images (DAE)}
\author{Alban Komla Henovi \and Asta Walo Diop \and Arthur Bomemo Walé}
\date{\today}

\begin{document}
\maketitle

\section*{Objectif de ce document}
\begin{itemize}
  \item Servir de \textbf{script} : quoi dire, dans quel ordre, et comment relier les slides.
  \item Clarifier les \textbf{notations} et \textbf{formules} (AE, DAE, MSE, PSNR, bruit).
  \item Donner des éléments de \textbf{pédagogie} : analogies, points d'attention, réponses aux questions.
\end{itemize}

\section*{Notations (à garder en tête pendant tout l'exposé)}
\begin{itemize}
  \item $x$ : image \textbf{propre} (référence, cible). Typiquement $x\in[0,1]^{H\times W\times C}$.
  \item $\tilde{x}$ : image \textbf{bruitée} (observation) obtenue en appliquant un bruit à $x$.
  \item $E_\theta$ : \textbf{encodeur} paramétré par $\theta$ (réduit/compresse l'information).
  \item $D_\theta$ : \textbf{décodeur} paramétré par $\theta$ (reconstruit une image).
  \item $z$ : représentation \textbf{latente} (vecteur / tenseur) dans le goulot d'étranglement.
  \item $\hat{x}$ : reconstruction / sortie du modèle.
  \item $\ell(\cdot,\cdot)$ : fonction de perte (ex. MSE). $\mathcal{L}(\theta)$ : perte moyenne sur le dataset.
\end{itemize}

\newpage
\tableofcontents
\newpage

% -----------------------------------------------------------------------------
\section{Slide de titre (logos + titre)}
\subsection*{Ce qu'il faut dire}
\begin{itemize}
  \item Présenter le thème : \emph{autoencodeurs} (apprentissage de représentations) et application au \emph{débruitage}.
  \item Situer le contexte : traitement d'images optiques, et objectif pédagogique (comprendre l'idée + voir une démo).
\end{itemize}

\subsection*{Transition}
\begin{itemize}
  \item \textbf{Transition} : « On commence par rappeler pourquoi on débruite, puis on introduit l'autoencodeur et sa variante débruiteuse. »
\end{itemize}

% -----------------------------------------------------------------------------
\section{Plan}
\subsection*{Ce qu'il faut dire}
\begin{itemize}
  \item Donner le fil conducteur : (1) contexte, (2) AE, (3) DAE, (4) métriques, (5) démo, (6) conclusion.
\end{itemize}

% -----------------------------------------------------------------------------
\section{Contexte et objectif}

\subsection{Pourquoi le débruitage en imagerie optique ?}
\paragraph{Message clé}
Le bruit est inévitable en acquisition réelle. On veut récupérer une image plus exploitable (visuellement et pour des tâches aval : segmentation, classification, détection).

\paragraph{Formule de la slide}
\[
  x_{\text{clean}} \approx f_\theta(\tilde{x})
\]
\begin{itemize}
  \item $f_\theta$ représente un modèle paramétré (ici un réseau profond) entraîné sur des exemples.
  \item $\tilde{x}$ est une observation bruitée, $x_{\text{clean}}$ est la cible souhaitée.
  \item « $\approx$ » car on ne récupère jamais parfaitement l'image propre : on fait une \textbf{estimation}.
\end{itemize}

\paragraph{Comment l'expliquer simplement}
\begin{itemize}
  \item Imaginez un filtre intelligent : au lieu d'un filtre fixe (médian, gaussien), on apprend un filtre adapté aux données.
  \item Le modèle apprend une transformation qui retire le bruit tout en préservant les structures importantes.
\end{itemize}

\subsection{Objectif de l'exposé}
\begin{itemize}
  \item Définir AE et DAE, expliquer l'intuition.
  \item Montrer une démo : ajout de bruit \textrightarrow{} entraînement \textrightarrow{} comparaison avant/après.
\end{itemize}

% -----------------------------------------------------------------------------
\section{Autoencodeurs : principe}

\subsection{Définition : autoencodeur}
\paragraph{Idée}
Un autoencodeur est un modèle qui apprend à \textbf{reconstruire} son entrée. S'il est contraint (bottleneck), il doit apprendre une représentation utile.

\paragraph{Formalisme (formules de la slide)}
\begin{align*}
  z &= E_\theta(x), \quad z \in \R^d \\
  \hat{x} &= D_\theta(z) \\
         &= D_\theta\big(E_\theta(x)\big)
\end{align*}
\begin{itemize}
  \item $E_\theta$ : compresse $x$ en $z$ (on réduit la dimension ou l'information).
  \item $D_\theta$ : décompresse $z$ et reconstruit $\hat{x}$.
  \item $d$ : dimension latente. Si $d$ est petit, on force un résumé (compression).
\end{itemize}

\paragraph{Questions fréquentes}
\begin{itemize}
  \item \textbf{Pourquoi "auto-supervisé" ?} Parce que la cible est l'entrée elle-même : pas besoin d'annotations.
  \item \textbf{Est-ce de la compression ?} Oui, en partie. Mais le but peut aussi être d'apprendre des features.
\end{itemize}

\subsection{Schéma AE}
\begin{itemize}
  \item Utiliser le schéma pour pointer : entrée $x$ \textrightarrow{} encodeur \textrightarrow{} latent $z$ \textrightarrow{} décodeur \textrightarrow{} sortie $\hat{x}$.
  \item \textbf{Message} : si le bottleneck est pertinent, le réseau retient l'essentiel (formes, textures) et ignore une partie du bruit.
\end{itemize}

\subsection{Perte de reconstruction}
\paragraph{Formule de la slide}
\[
  \mathcal{L}(\theta)=\frac{1}{m}\sum_{j=1}^m \ell\big(x^{(j)}, \hat{x}^{(j)}\big)
\]
\begin{itemize}
  \item $m$ : nombre d'exemples.
  \item $x^{(j)}$ : exemple propre, $\hat{x}^{(j)}$ : reconstruction correspondante.
  \item On minimise $\mathcal{L}$ en ajustant les paramètres $\theta$ (descente de gradient, Adam, etc.).
\end{itemize}

\paragraph{Exemple : MSE}
\[
  \ell(x,\hat{x}) = \frac{1}{N}\norm{x-\hat{x}}_2^2
\]
\begin{itemize}
  \item $N$ : nombre total de pixels (fois canaux) dans l'image.
  \item $\norm{\cdot}_2^2$ : somme des carrés des différences pixel à pixel.
  \item \textbf{Interprétation} : pénalise fortement les grosses erreurs.
\end{itemize}

\paragraph{Limite importante}
La MSE est simple mais peut produire des sorties "trop lisses" (moyennage) car elle encourage la moyenne des solutions plausibles.

% -----------------------------------------------------------------------------
\section{Autoencodeur débruiteur (DAE)}

\subsection{Idée : Denoising Autoencoder}
\paragraph{Intuition}
Au lieu de reconstruire $x$ à partir de $x$, on reconstruit $x$ à partir d'une version corrompue $\tilde{x}$. Cela force le modèle à apprendre à enlever le bruit.

\paragraph{Formule de la slide}
\[
  \min_\theta\; \E_{x\sim p_{\text{data}}}\; \E_{\tilde{x}\sim q(\cdot\mid x)}\; \ell\big(x,\, D_\theta(E_\theta(\tilde{x}))\big)
\]
\begin{itemize}
  \item $p_{\text{data}}$ : distribution des données propres (les images qu'on veut).
  \item $q(\tilde{x}\mid x)$ : "mécanisme" de corruption : comment on fabrique $\tilde{x}$ à partir de $x$.
  \item Double espérance : on moyenne sur les images et sur les réalisations de bruit.
  \item $D_\theta(E_\theta(\tilde{x}))$ : sortie du réseau quand l'entrée est bruitée.
\end{itemize}

\paragraph{Comment le dire à l'oral}
\begin{itemize}
  \item « On simule du bruit, et on apprend une fonction qui inverse ce bruit au mieux. »
  \item « Le réseau voit une image dégradée mais il est pénalisé par rapport à l'image propre. »
\end{itemize}

\subsection{Schéma DAE}
\begin{itemize}
  \item Insister sur la différence : l'entrée est $\tilde{x}$, la cible est $x$.
  \item Point d'attention : si le bruit d'entraînement n'a rien à voir avec le bruit réel, le modèle généralise mal.
\end{itemize}

\subsection{Modèles de bruit (exemples)}
\paragraph{Bruit gaussien additif}
\[
  \tilde{x} = x + \sigma\,\varepsilon,\quad \varepsilon\sim\mathcal{N}(0,I)
\]
\begin{itemize}
  \item $\varepsilon$ : bruit aléatoire de moyenne nulle.
  \item $\sigma$ : contrôle l'intensité (plus $\sigma$ est grand, plus l'image est bruitée).
  \item Après ajout, on \textbf{clip} souvent dans $[0,1]$ pour rester dans une dynamique d'image valide.
\end{itemize}

\paragraph{Sel et poivre}
\begin{itemize}
  \item Certains pixels sont remplacés par 0 ou 1 (noir/blanc) avec une certaine probabilité.
\end{itemize}

\paragraph{Bruit multiplicatif (speckle)}
\[
  \tilde{x} = x\odot(1+\varepsilon)
\]
\begin{itemize}
  \item "$\odot$" : multiplication terme à terme.
  \item Le bruit dépend de l'intensité du signal (plus réaliste dans certains capteurs).
\end{itemize}

% -----------------------------------------------------------------------------
\section{Architecture : AE convolutionnel}

\subsection{Pourquoi des convolutions ?}
\begin{itemize}
  \item En image, les motifs locaux (bords, textures) comptent : les convolutions les capturent bien.
  \item Partage de poids : beaucoup moins de paramètres qu'un MLP à taille d'image identique.
  \item Encodeur : réduit la résolution (pooling/strides) ; décodeur : restaure la résolution (upsampling / transpose conv).
\end{itemize}

\subsection{Pseudo-code (entraînement DAE)}
\paragraph{Ce qu'il illustre}
\begin{itemize}
  \item Pour chaque batch d'images propres $x$, on fabrique $x_{\text{tilde}}$ bruitée.
  \item Le modèle prédit $x_{\hat{}}$.
  \item On calcule la MSE entre $x_{\hat{}}$ et $x$ (cible = propre), puis on met à jour les poids.
\end{itemize}

\paragraph{Point pédagogique}
\begin{itemize}
  \item Dire clairement : « entrée = bruitée, cible = propre » (c'est la clé du DAE).
  \item Mentionner la généralisation : on peut tester à des niveaux de bruit un peu différents.
\end{itemize}

% -----------------------------------------------------------------------------
\section{Évaluation et métriques}

\subsection{MSE}
\begin{itemize}
  \item La MSE est une moyenne des erreurs au carré : plus elle est faible, plus la reconstruction est proche.
  \item Avantage : simple, dérivable, standard.
  \item Limite : pas toujours alignée avec la perception humaine.
\end{itemize}

\subsection{PSNR (Peak Signal-to-Noise Ratio)}
\paragraph{Formule (slide)}
Si l'image est dans $[0,1]$ et $MAX_I=1$ :
\[
\mathrm{PSNR}(x,\hat{x}) = 10\log_{10}\Big(\frac{MAX_I^2}{\mathrm{MSE}(x,\hat{x})}\Big)
\]
\paragraph{Explication}
\begin{itemize}
  \item Le PSNR exprime la qualité en \textbf{décibels} (échelle logarithmique).
  \item Quand la MSE diminue, le PSNR augmente.
  \item Interprétation : c'est une façon standard de comparer des méthodes de débruitage/compression.
\end{itemize}

\paragraph{Intuition rapide}
\begin{itemize}
  \item Si l'erreur est très petite, $\frac{1}{\mathrm{MSE}}$ est grand \textrightarrow{} log grand \textrightarrow{} PSNR élevé.
\end{itemize}

\subsection{SSIM}
\begin{itemize}
  \item SSIM compare des structures locales (contraste, luminance, structure).
  \item Souvent mieux corrélé à la perception, mais plus complexe à expliquer/calculer.
  \item Dans la démo, on se concentre sur MSE/PSNR pour rester simple.
\end{itemize}

% -----------------------------------------------------------------------------
\section{Démo notebook}

\subsection{Démo : déroulé}
\paragraph{Ce qu'il faut raconter (pas à pas)}
\begin{enumerate}
  \item \textbf{Chargement} : on prend une image locale (\texttt{assets/inputs/image1.webp}) et on normalise dans $[0,1]$.
  \item \textbf{Bruit} : on ajoute un bruit gaussien (contrôlé par $\sigma$), puis on clip.
  \item \textbf{Patchs} : on extrait des patchs bruité/propre alignés pour former un dataset (utile si on n'a qu'une image).
  \item \textbf{Entraînement} : AE convolutionnel, perte MSE (entrée bruitée, cible propre).
  \item \textbf{Inférence} : on applique le modèle à l'image entière et on compare les 3 versions.
\end{enumerate}

\paragraph{Pourquoi les patchs ?}
\begin{itemize}
  \item Avec une seule grande image, on peut créer beaucoup d'exemples en découpant en patchs.
  \item Cela rend l'entraînement possible sans gros dataset.
  \item En contrepartie, il faut faire attention aux effets de bord et au "style" unique de l'image.
\end{itemize}

\subsection{Résultats (3 images)}
\begin{itemize}
  \item \textbf{Entrée bruitée} : montrer que la structure existe encore mais est dégradée.
  \item \textbf{Sortie AE} : expliquer qu'on retire du bruit mais qu'il peut y avoir du lissage.
  \item \textbf{Cible propre} : rappeler que c'est la référence (ground truth) dans notre expérience (bruit simulé).
\end{itemize}

\paragraph{Questions fréquentes}
\begin{itemize}
  \item \textbf{Pourquoi on sait la cible propre ?} Parce qu'on ajoute nous-mêmes le bruit à une image propre.
  \item \textbf{Et si on n'a pas d'image propre ?} Il existe des méthodes "blind" / Noise2Noise, etc. (hors scope de cet exposé).
\end{itemize}

% -----------------------------------------------------------------------------
\section{Conclusion}
\begin{itemize}
  \item Reformuler la différence AE vs DAE (entrée = propre vs entrée = bruitée, cible = propre).
  \item Souligner que la qualité dépend du bruit simulé, de l'architecture, et des données.
\end{itemize}

\section{Références (liens)}
\begin{itemize}
  \item Ichi.pro (FR) : \url{https://ichi.pro/fr/autoencodeurs-de-debruitage-tensorflow-102261356342629}
  \item LaBRI (P. Preuter) : \url{https://www.labri.fr/perso/preuter/dl2025/cours4-dl.html}
  \item Colab / Eyabennessib : \url{https://colab.research.google.com/github/Eyabennessib/Machine-Learning/blob/main/TP7_autoencodeurs_EBN.ipynb/}
  \item Simon Thomine : \url{https://simonthomine.github.io/CoursDeepLearning/fr/04_Autoencodeurs/02_DenoisingAE.html}
  \item DataCamp (FR) : \url{https://www.datacamp.com/fr/tutorial/introduction-to-autoencoders}
  \item NYU DLSP20 (FR) : \url{https://atcold.github.io/NYU-DLSP20/fr/week07/07-3/}
\end{itemize}

\section*{Fin (Merci / Questions)}
\begin{itemize}
  \item Slide "Merci" : laisser 2--3 secondes, puis annoncer la session de questions.
  \item Slide "Questions ?" : proposer 2 questions d'amorçage :
  \begin{itemize}
    \item « Quel type de bruit rencontre-t-on le plus en imagerie optique ? »
    \item « Qu'est-ce qu'on perd quand on optimise uniquement la MSE ? »
  \end{itemize}
\end{itemize}

\end{document}
